% ID: SEL-D0081-OB-E04
% Tanda S081.
\textbf{a)} La población total es $400+320=720$. La afijación proporcional distribuye el tamaño de la muestra según el peso de cada estrato:
\[
 n_H=54\frac{400}{720}=30,\qquad
 n_M=54\frac{320}{720}=24.
\]
Se seleccionan aleatoriamente 30 hombres y 24 mujeres dentro de sus respectivos estratos.

\textbf{b)} El enunciado no especifica reposición ni orden (INC-027). Adoptamos muestreo aleatorio simple sin reemplazo, con subconjuntos equiprobables. El orden no afecta a la media y las seis muestras son
\[
 \begin{array}{c|cccccc}
 \text{Muestra}&\{1,2\}&\{1,3\}&\{1,4\}&\{2,3\}&\{2,4\}&\{3,4\}\\
 \bar X&1{,}5&2&2{,}5&2{,}5&3&3{,}5
 \end{array}
\]
Su esperanza es $2{,}5$ y su varianza es
\[
 \operatorname{Var}(\bar X)
 =\frac{(1{,}5-2{,}5)^2+(2-2{,}5)^2+0+0+(3-2{,}5)^2+(3{,}5-2{,}5)^2}{6}
 =\boxed{\frac5{12}}.
\]
Si se considera el orden sin reemplazo, aparecen doce pares, dos por subconjunto, y se obtiene la misma varianza.

\textit{Convención alternativa con reemplazo.} En ese caso hay dieciséis pares ordenados equiprobables, $(i,j)$ con $i,j\in\{1,2,3,4\}$. La tabla enumera todos los pares y sus medias mediante la fila $i$ y la columna $j$:
\[
 \begin{array}{c|cccc}
 i\backslash j&1&2&3&4\\\hline
 1&1&1{,}5&2&2{,}5\\
 2&1{,}5&2&2{,}5&3\\
 3&2&2{,}5&3&3{,}5\\
 4&2{,}5&3&3{,}5&4
 \end{array}
\]
La varianza de la población es $5/4$. Con dos extracciones independientes, $\operatorname{Var}(\bar X)=(5/4)/2=\boxed{5/8}$. Esta alternativa requiere el supuesto de reposición.
